\documentclass[11pt,a4paper]{article}

\PassOptionsToPackage{warnings-off={mathtools-colon,mathtools-overbracket}}{unicode-math}
\usepackage{mathtools}
\usepackage{newcomputermodern}
\usepackage[a4paper,margin=26mm]{geometry}
\usepackage{microtype}
\usepackage{xcolor}
\usepackage{graphicx}
\usepackage{booktabs}
\usepackage{siunitx}
\usepackage{csquotes}
\usepackage[backend=biber,style=authoryear,maxcitenames=2]{biblatex}
\usepackage{hyperref}
\usepackage{bookmark}

\addbibresource{references.bib}
\graphicspath{{figures/}}

\definecolor{Accent}{HTML}{285EA8}
\hypersetup{
  colorlinks=true,
  linkcolor=Accent,
  urlcolor=Accent,
  citecolor=Accent,
  pdftitle={A Clear and Specific Article Title},
  pdfauthor={First Author and Second Author}
}

\setlength{\parindent}{1.2em}
\setlength{\parskip}{0.2em}

\newcommand{\keywords}[1]{%
  \par\smallskip
  \noindent\textbf{Keywords:} #1
}

\title{A Clear and Specific Article Title}
\author{
  First Author\\
  \small Department, University\\
  \small \href{mailto:first.author@example.com}{first.author@example.com}
  \and
  Second Author\\
  \small Department, University
}
\date{\today}

\begin{document}

\maketitle

\begin{abstract}
State the research question, method, principal result, and main implication in
one compact paragraph. The abstract should stand on its own and should not
introduce citations or undefined abbreviations.
\keywords{LaTeX, scientific writing, reproducible documents}
\end{abstract}

\section{Introduction}

Explain the problem, establish why it matters, and identify the gap addressed
by the study. End the introduction with a precise objective or research
question. Structured authoring helps separate document content from its
presentation \parencite{lamport1994latex}.

\section{Methods}

Describe the data, apparatus, participants, or procedure in enough detail for
the work to be evaluated and reproduced. Define the primary quantity before
using it. For example, a normalized change can be written as
\begin{equation}
  r = \frac{x_{\mathrm{after}}-x_{\mathrm{before}}}{x_{\mathrm{before}}}.
  \label{eq:normalized-change}
\end{equation}

State inclusion rules, uncertainty treatment, and analysis decisions here
rather than after the results are known.

\section{Results}

Present observations before interpretation. Every table and figure should be
referenced from the text and should communicate one main point.

\begin{table}[htbp]
  \centering
  \caption{Illustrative measurements for the article template.}
  \label{tab:results}
  \begin{tabular}{@{}l S[table-format=2.1] S[table-format=2.1] S[table-format=1.2]@{}}
    \toprule
    Group & {Before} & {After} & {Change} \\
    \midrule
    A & 12.4 & 15.1 & 0.22 \\
    B & 11.8 & 13.0 & 0.10 \\
    C & 12.1 & 12.5 & 0.03 \\
    \bottomrule
  \end{tabular}
\end{table}

The illustrative values in \autoref{tab:results} show how aligned numeric
columns make comparisons easier. Replace them with measured values and report
sample sizes and uncertainty where appropriate.

\section{Discussion}

Interpret the results in relation to the research question. Distinguish what
the evidence directly supports from plausible explanations, compare the
result with relevant literature, and identify limitations that affect the
conclusion.

\section{Conclusion}

Answer the research question directly and state the practical or theoretical
implication. Do not introduce new evidence in the conclusion.

\printbibliography

\end{document}
